\documentclass{beamer}

\usepackage{../tema}


\title{SAL}
\author[L. Granziero]{L. Granziero}
\setDate{04/03/2026}

\begin{document}

{
  \setbeamertemplate{footline}{}
  \begin{frame}
    \titlepage
  \end{frame}
}

\addMember{Filippo}{Venzo}{Progettista}
\addMember{Giuseppe}{De Fina}{Progettista}
\addMember{Valeria}{Baleanu}{Progettista}
\addMember{Christian}{Libralato}{Verificatore}
\addMember{Leonardo}{Pellizzon}{Amministratore}
\addMember{Francesco}{Pasqual}{Progettista}
\addMember{Luca}{Granziero}{Responsabile}
\makeMembersFrame

\begin{frame}{Avanzamenti settimanali}
    \begin{itemize}
        \item Continuazione fase di progettazione, sempre con il coinvolgimento di 4 progettisti 
\end{itemize}
\end{frame}

% PRIMO FRAME: Attività terminate dal verbale esterno 

% Reset delle entries per il secondo frame 
\renewcommand{\precedentTodoEntries}{}

%  SECONDO FRAME: Attività terminate dal verbale interno
\addPrecedentTodo{vi\_2026\_02\_25.a1}{\#141}{Scrittura verbale interno}{L. Granziero}{28/02/2026}

\makePrecedentTodoFrame


\begin{frame}{Attività da completare}
\begin{itemize}
    \item Aggiornare il PdP allo sprint 8
    \item Disporre di una progettazione a un livello di dettaglio adeguato per poter avviare la programmazione all’inizio del prossimo sprint (o al più tardi a metà sprint).
\end{itemize}
\end{frame}

\begin{frame}{Dubbi incontrati}
  \begin{itemize}
    \item dubbi sulla sezione analytics e consigli 
    \item come vengono memorizzate le "anomalie di impianto" da visualizzare nelle analytics? 
    \item Per ogni tipologia di analytics abbiamo associato un livello di visualizzazione (dispositivo, impianto o reparto). Preferite mantenere questa associazione fissa oppure prevedere la possibilità di scegliere il livello per ciascuna analytics?
  \end{itemize}
\end{frame}

\end{document}